% =====================================================================
%  Bien dich: pdfLaTeX (mac dinh cua Overleaf) -- khong can chinh gi.
%  Neu muon dung font he thong thi doi compiler sang XeLaTeX va thay
%  ba dong inputenc/fontenc/lmodern ben duoi bang:
%      \usepackage{fontspec}
%      \setmainfont{TeX Gyre Termes}
%
%  Moi khoi verbatim trong file nay deu la ASCII thuan y, vi verbatim
%  khong xu ly duoc ky tu tieng Viet nhieu byte.
% =====================================================================
\documentclass[11pt,a4paper]{article}

\usepackage[utf8]{inputenc}
\usepackage[T5]{fontenc}
\usepackage{lmodern}

\usepackage[margin=2.5cm]{geometry}
\usepackage{amsmath}
\usepackage{amssymb}
\usepackage{booktabs}
\usepackage{array}
\usepackage{longtable}
\usepackage{caption}
\usepackage{microtype}
\usepackage{enumitem}
\usepackage{framed}
\usepackage{tikz}
\usetikzlibrary{arrows.meta, positioning}
\usepackage[hidelinks]{hyperref}

\renewcommand{\tablename}{Bảng}
\renewcommand{\figurename}{Hình}
\renewcommand{\contentsname}{Nội dung}
\renewcommand{\appendixname}{Phụ lục}

\captionsetup[table]{skip=6pt}
\captionsetup[figure]{skip=8pt}
\setlist[itemize]{itemsep=3pt, topsep=4pt}
\setlist[enumerate]{itemsep=4pt, topsep=4pt}

\newcolumntype{R}{>{\raggedleft\arraybackslash}p{2.7cm}}
\newcolumntype{L}{>{\raggedright\arraybackslash}p{7.4cm}}

% Khoi luu y, dung leftbar cua goi framed cho nhe
\newenvironment{luuy}{\begin{leftbar}\noindent\ignorespaces}{\end{leftbar}}

\title{\bfseries Hệ thống này làm gì\\[2pt]
\large\mdseries Giải thích từ đầu cho người chưa biết về đề tài}
\author{Nguyễn Thiện Trường}
\date{Tháng 8, 2026}

\begin{document}
\maketitle

\noindent
Tài liệu này viết cho người chưa biết gì về đề tài. Đọc hết sẽ hiểu: bài toán là
gì, bài báo gốc HRM-PET hoạt động ra sao, chúng ta phát hiện điều gì, thêm cái gì
vào, thêm ở chỗ nào trong mã nguồn, và tại sao lại thêm đúng chỗ đó. Không cần
đọc mã nguồn để hiểu tài liệu, nhưng mọi chỗ nhắc tới mã đều kèm tên file và số
dòng để tra khi cần.

\tableofcontents
\bigskip

% =====================================================================
\section{Bài toán: học liên tục}

\subsection{Vấn đề}

Một mô hình nhận diện ảnh bình thường được huấn luyện một lần trên toàn bộ dữ
liệu rồi đem dùng. \textbf{Học liên tục} (continual learning) là tình huống khác:
dữ liệu đến theo từng đợt, và mô hình phải học đợt mới mà không được xem lại đợt
cũ.

Ví dụ với bộ Split-CIFAR100 mà chúng ta dùng: tác vụ 1 học lớp 0--9, tác vụ 2 học
lớp 10--19, cứ thế cho tới tác vụ 10 học lớp 90--99. Khi học tác vụ 2, mô hình
\textbf{không còn được xem ảnh của lớp 0--9 nữa}. Đến cuối cùng nó vẫn phải phân
loại đúng cả 100 lớp.

\subsection{Vì sao khó: quên thảm khốc}

Nếu cứ huấn luyện tiếp một cách ngây thơ, mạng nơ-ron sẽ ghi đè trọng số cũ để
tối ưu cho dữ liệu mới. Kết quả là nó học tác vụ 10 rất tốt nhưng gần như quên
sạch tác vụ 1. Hiện tượng này gọi là \textbf{quên thảm khốc} (catastrophic
forgetting) và là trở ngại trung tâm của cả lĩnh vực.

\subsection{Ràng buộc: không lưu mẫu}

Cách chống quên đơn giản nhất là giữ lại vài trăm ảnh cũ để ôn lại, gọi là
\emph{replay}. Nhưng nhiều tình huống thực tế cấm điều đó: dữ liệu y tế không
được lưu vì lý do riêng tư, dữ liệu người dùng bị giới hạn bởi luật.

Đề tài này theo giao thức \textbf{không lưu mẫu} (exemplar-free) nghiêm ngặt:
không giữ bất kỳ ảnh nào, cũng không giữ đặc trưng của từng ảnh riêng lẻ. Chỉ
được giữ các đại lượng thống kê tổng hợp, ví dụ trung bình và hiệp phương sai của
cả một lớp. Toàn bộ mã đánh giá chạy dưới cờ
\texttt{-{}-strict\_exemplar\_free}, và có một bước kiểm tra checkpoint sẽ báo
lỗi nếu phát hiện đặc trưng theo mẫu bị lưu lén.

\subsection{Cách tiếp cận hiện đại: đóng băng mạng lớn, gắn bộ điều hợp nhỏ}

Thay vì huấn luyện lại toàn bộ mạng, hướng hiện đại làm như sau:

\begin{enumerate}
\item Lấy một mạng thị giác lớn đã huấn luyện sẵn trên tập dữ liệu khổng lồ --
      ở đây là \textbf{ViT-B/16} huấn luyện trên ImageNet-21K, khoảng 14 triệu ảnh.
\item \textbf{Đóng băng} toàn bộ mạng đó, không đụng vào một trọng số nào.
\item Với mỗi tác vụ mới, huấn luyện một \textbf{bộ điều hợp} (adapter) rất nhỏ
      gắn thêm vào. Ở đây dùng \textbf{LoRA hạng 8}, chỉ vài trăm nghìn tham số.
\end{enumerate}

Vì mạng gốc không bao giờ thay đổi, kiến thức cũ trong nó không thể bị ghi đè.
Mỗi tác vụ có LoRA riêng, và LoRA của tác vụ 1 vẫn nguyên vẹn sau khi học tác vụ
10. Nhưng cách này đẻ ra một vấn đề mới, và đó chính là chỗ HRM-PET bước vào.

% =====================================================================
\section{HRM-PET hoạt động như thế nào}

\subsection{Vấn đề mà HRM-PET giải quyết}

Ta có 10 bộ LoRA, mỗi bộ chuyên cho một tác vụ. Khi một ảnh mới đến lúc kiểm tra,
\textbf{ta không được cho biết ảnh đó thuộc tác vụ nào}, phải tự đoán để chọn
đúng LoRA.

Hình dung như một thư viện có 10 chuyên gia, mỗi người giỏi một lĩnh vực. Có
người mang ảnh đến hỏi nhưng không nói ảnh thuộc lĩnh vực nào. Người trực cửa
phải đoán để chuyển đến đúng chuyên gia; đoán sai thì câu trả lời có nguy cơ sai.
Bước đoán đó gọi là \textbf{suy luận danh tính tác vụ} (task-identity inference,
viết tắt TII).

\subsection{Bốn bước}

\paragraph{Bước 1 -- TII: đoán tác vụ.}
Một mô hình phụ (trong mã gọi là \texttt{original\_model}) chạy ảnh qua và đo
\emph{năng lượng prompt}. Ý tưởng: mỗi tác vụ có một tập prompt riêng, và prompt
của tác vụ đúng sẽ phản ứng mạnh hơn với ảnh thuộc tác vụ đó. Hàm đo dùng ở đây
là entropy suy rộng (generalized entropy, cờ \texttt{-{}-En gen}).

\paragraph{Bước 2 -- Định tuyến.}
Chọn tác vụ có điểm cao nhất, nạp LoRA của tác vụ đó vào mạng, rồi chạy ảnh qua.

\paragraph{Bước 3 -- Re-matching bằng DRM và CRM.}
Đây là đóng góp chính của bài báo gốc. Sau khi chạy LoRA đầu tiên, mô hình đưa ra
một lớp dự đoán. Mỗi lớp thuộc về đúng một tác vụ, nên \textbf{bản thân lớp được
dự đoán cũng là một phiếu bầu cho tác vụ}.

\begin{itemize}
\item \textbf{DRM} (Direct Re-Matching) dùng luôn lớp dự đoán đó để định tuyến
      lại: nếu LoRA của tác vụ 3 lại dự đoán ra một lớp thuộc tác vụ 7, thì có lẽ
      nên thử LoRA của tác vụ 7.
\item \textbf{CRM} (Confidence Re-Matching) chỉ kích hoạt việc định tuyến lại khi
      độ tin cậy thấp, đo bằng chính hàm GEN. Dự đoán đã chắc chắn thì để yên.
\end{itemize}

Đây là lý do HRM-PET tốn trung bình \textbf{1.83 lượt gọi LoRA cho mỗi ảnh} trên
ImageNet-R chứ không phải 1.0: một số ảnh phải chạy lại lần hai.

\paragraph{Bước 4 -- Phân lớp.}
Đưa đặc trưng cuối vào một \textbf{đầu phân lớp dùng chung}, trải trên
\emph{toàn bộ} các lớp đã gặp từ đầu đến giờ.

\begin{luuy}
\textbf{Ghi nhớ chi tiết này.} Đầu phân lớp là \textbf{dùng chung cho mọi tác
vụ}, không phải mỗi tác vụ một đầu riêng. Toàn bộ Mục \ref{sec:phat-hien-2} xoay
quanh hệ quả của nó.
\end{luuy}

\begin{figure}[htbp]
\centering
\begin{tikzpicture}[
  font=\small, >={Stealth[length=2mm]},
  box/.style={draw, minimum height=8mm, minimum width=3.4cm, align=center},
  lbl/.style={font=\footnotesize, align=left}
]
\node[box] (tii) at (0,0)    {TII: đoán tác vụ};
\node[box] (lo)  at (0,-1.5) {Chọn LoRA};
\node[box] (drm) at (0,-3.0) {DRM / CRM};
\node[box] (cl)  at (0,-4.5) {Đầu phân lớp chung};
\node[box] (out) at (0,-6.0) {Dự đoán};

\draw[->] (tii) -- (lo);
\draw[->] (lo)  -- (drm);
\draw[->] (drm) -- (cl);
\draw[->] (cl)  -- (out);

\node[lbl, right=6mm of tii] {đo năng lượng prompt};
\node[lbl, right=6mm of lo]  {nạp bộ điều hợp của tác vụ đó};
\node[lbl, right=6mm of drm] {dùng lớp dự đoán để định tuyến lại\\nếu chưa chắc chắn};
\node[lbl, right=6mm of cl]  {trải trên MỌI lớp đã gặp};

\node[left=8mm of tii] (in) {ảnh vào};
\draw[->] (in) -- (tii);
\end{tikzpicture}
\caption{Quy trình gốc của HRM-PET.}
\label{fig:goc}
\end{figure}

% =====================================================================
\section{Ba phát hiện dẫn tới đề xuất}

\subsection{Phát hiện 1 -- TII không phải tín hiệu định tuyến tốt nhất}

Chúng ta thử một cách đoán tác vụ hoàn toàn khác, gọi là \textbf{đầu RP} (mô tả
kỹ ở Mục \ref{sec:dau-rp}). Nó không dùng prompt, không dùng năng lượng, mà dùng
đặc trưng ảnh và một phép hồi quy tuyến tính.

\begin{table}[htbp]
\centering
\caption{Độ chính xác đoán tác vụ (\%), seed 42.}
\label{tab:routing}
\begin{tabular}{lrrr}
\toprule
Cách đoán tác vụ & ImageNet-R & CIFAR-100 & CUB-200 \\
\midrule
TII (bài báo gốc, bước 1)  & 63.80 & 81.21 & 92.96 \\
HRM-PET sau khi có DRM/CRM & 77.79 & 89.82 & 93.21 \\
Đầu RP                     & 77.12 & 89.51 & 94.02 \\
\midrule
Hợp của TII và đầu RP      & \textbf{82.99} & \textbf{92.86} & \textbf{96.26} \\
\bottomrule
\end{tabular}
\end{table}

Hàng cuối cần giải thích. Đó \textbf{không phải} một phương pháp, mà là phép đo
``nếu có một vị thần biết chọn giữa hai câu trả lời thì đúng được bao nhiêu'':
tỉ lệ ảnh mà \emph{ít nhất một} trong hai cách đoán ra kết quả đúng.

Con số 82.99 so với 77.79 nói lên điều then chốt: \textbf{hai cách đoán sai ở
những ảnh khác nhau}. Trên ImageNet-R, đầu RP đoán đúng $19.2\,\%$ số ảnh mà TII
đoán sai. Đây là điều kiện bắt buộc để việc kết hợp có ý nghĩa; nếu hai nguồn
cùng sai ở cùng những ảnh thì gộp lại chẳng thêm được gì.

\subsection{Phát hiện 2 -- định tuyến tốt hơn không tự động thành phân lớp tốt hơn}
\label{sec:phat-hien-2}

Sau khi cải thiện định tuyến, chúng ta gặp một điều bất ngờ: \textbf{định tuyến
tăng hơn 2 điểm, nhưng độ chính xác phân lớp chỉ tăng 0.29 điểm.} Trên CUB-200
còn tệ hơn: định tuyến tăng 1.13 điểm mà phân lớp \emph{giảm} 0.05 điểm.

Nguyên nhân nằm ở chi tiết đã đánh dấu: \textbf{đầu phân lớp dùng chung}.

Giả sử một ảnh thuộc tác vụ 3 nhưng bị định tuyến nhầm sang LoRA của tác vụ 7.
LoRA số 7 vẫn là một bộ điều hợp nhỏ gắn trên cùng mạng xương sống, nên đặc trưng
nó tạo ra không quá khác biệt. Đầu phân lớp lại nhìn thấy \emph{tất cả} 200 lớp,
nên nó hoàn toàn có thể vẫn chọn đúng lớp thuộc tác vụ 3. Từ đó suy ra hai điều:

\begin{itemize}
\item Sửa đường định tuyến cho ảnh đó \textbf{không mua được gì}, nó vốn đã đúng.
\item Ngược lại, định tuyến lại một ảnh vốn đã đúng \textbf{có thể làm nó sai đi}.
\end{itemize}

Đó là lý do phần lợi ích bị thất thoát. Muốn độ chính xác phân lớp tăng thật thì
phải can thiệp vào chính bước phân lớp, không chỉ bước định tuyến.

\subsection{Phát hiện 3 -- đầu RP còn là một bộ phân lớp bị bỏ phí}

Đầu RP không chỉ đoán tác vụ. Bản thân nó là \textbf{một bộ phân lớp hoàn chỉnh},
cho điểm số cho từng lớp cụ thể chứ không chỉ từng tác vụ. Chúng ta đang dùng nó
để đoán tác vụ rồi vứt phần còn lại đi.

\begin{table}[htbp]
\centering
\caption{Mức bổ trợ ở cấp phân lớp (\%), seed 42. Cột cuối là tỉ lệ ảnh mà đầu
phân lớp chính sai còn đầu RP đúng.}
\label{tab:cls}
\begin{tabular}{lrrrr}
\toprule
Bộ dữ liệu & Đầu phân lớp chính & Đầu RP & Hợp & Đầu RP cứu riêng \\
\midrule
ImageNet-R & 74.31 & 70.08 & 78.89 & $+4.58$ \\
CIFAR-100  & 89.97 & 88.76 & 92.40 & $+2.43$ \\
CUB-200    & 86.48 & \textbf{87.50} & 90.35 & $+3.87$ \\
\bottomrule
\end{tabular}
\end{table}

Điểm đáng chú ý: trên ImageNet-R đầu RP yếu hơn tới 4 điểm về tổng thể (70.08 so
với 74.31), \textbf{nhưng nó vẫn đúng ở $4.58\,\%$ số ảnh mà đầu chính bỏ lỡ}.
Một bộ phân lớp yếu hơn vẫn có giá trị nếu nó sai ở chỗ khác. Trên CUB-200 thì
đầu RP còn mạnh hơn cả đầu chính.

% =====================================================================
\section{Thành phần thứ nhất: đầu RP}
\label{sec:dau-rp}

\subsection{Ý tưởng}

Đầu RP là phiên bản của kỹ thuật trong bài báo \textbf{RanPAC}. Ý tưởng nghe hơi
phản trực giác nhưng rất hiệu quả: chiếu đặc trưng lên một không gian ngẫu nhiên
có số chiều rất lớn, rồi giải một bài toán hồi quy tuyến tính trong không gian
đó.

\begin{align}
h &= \mathrm{ReLU}(f W), &&\text{$W$ đóng băng, sinh lại từ seed}\\
G &= \sum_n h_n h_n^{\top}, \qquad
C = \sum_n h_n\,\mathrm{onehot}(y_n)^{\top},
&&\text{cộng dồn qua mọi ảnh đã thấy}\\
s &= h^{\top} (G + \lambda I)^{-1} C
&&\text{điểm số cho từng lớp}
\end{align}

Cụ thể ba bước. \emph{Một}, nhân đặc trưng 768 chiều với ma trận ngẫu nhiên $W$
để ra 10\,000 chiều rồi cho qua ReLU; $W$ được sinh một lần từ seed cố định và
\textbf{không bao giờ được huấn luyện}. \emph{Hai}, với mỗi ảnh huấn luyện, cộng
dồn vào ma trận Gram $G$ (ghi lại các chiều nào hay xuất hiện cùng nhau) và ma
trận tương quan lớp $C$ (ghi lại chiều nào ứng với lớp nào). \emph{Ba}, trọng số
phân lớp là nghiệm hồi quy ridge $(G+\lambda I)^{-1}C$.

\subsection{Vì sao cách này hợp với học liên tục}

\paragraph{Nó tuân thủ ràng buộc không lưu mẫu.}
$G$ và $C$ là \emph{tổng cộng dồn}. Khi tác vụ mới đến, chỉ cần cộng thêm vào,
không cần xem lại dữ liệu cũ. Ma trận $W$ sinh lại được từ seed nên cũng không
cần lưu. Không có ảnh nào, không có đặc trưng riêng lẻ nào được giữ.

\paragraph{Nó không hề quên.}
Không có bước huấn luyện lặp nào để ghi đè kiến thức cũ. Nghiệm hồi quy trên $G$
và $C$ đầy đủ tương đương với việc huấn luyện tuyến tính trên toàn bộ dữ liệu
cùng lúc.

\paragraph{Nó không tốn gì thêm.}
Toàn bộ phép tính là cộng ma trận và một lần giải hệ tuyến tính ở cuối mỗi tác vụ.

\subsection{Một chi tiết dễ sai}

Đặc trưng $f$ phải lấy từ \textbf{một bộ LoRA cố định duy nhất} (cờ
\texttt{-{}-rp\_lora\_task 0}, luôn dùng LoRA của tác vụ đầu tiên), chứ không
phải mỗi tác vụ dùng LoRA riêng của nó.

Lý do: $G$ và $C$ cộng dồn qua mọi tác vụ. Nếu mỗi tác vụ dùng một LoRA khác, các
đặc trưng sẽ nằm ở những không gian khác nhau và việc cộng dồn trở nên vô nghĩa.
Dùng một LoRA cố định chính là vai trò \textbf{thích nghi phiên đầu}
(first-session adaptation) trong RanPAC, chỉ khác là ở đây nó có sẵn từ HRM-PET
nên không phải huấn luyện thêm.

\subsection{Thêm vào đâu trong mã nguồn}

\begin{table}[htbp]
\centering
\caption{Vị trí của đầu RP trong mã nguồn.}
\begin{tabular}{lL}
\toprule
Việc & Vị trí \\
\midrule
Toàn bộ đầu RP        & \texttt{engines/random\_projection\_head.py} (file mới) \\
Sinh ma trận chiếu    & \texttt{\_ensure\_projection}, dòng 72 \\
Cộng dồn $G$ và $C$   & \texttt{accumulate\_rp\_statistics}, dòng 115 \\
Giải hồi quy ridge    & \texttt{solve\_rp\_head}, dòng 135 \\
Gọi ở cuối mỗi tác vụ & \texttt{trainers/lora\_trainer.py}, dòng 312--321 \\
Trích đặc trưng       & \texttt{pin\_rp\_extractor} / \texttt{rp\_extractor},
                        engine dòng 690 và 711 \\
\bottomrule
\end{tabular}
\end{table}

\noindent Cờ bật:
\texttt{-{}-rp\_head -{}-rp\_dim 10000 -{}-rp\_lambda 1e4}
\texttt{-{}-rp\_feature\_source lora -{}-rp\_lora\_task 0}

% =====================================================================
\section{Thành phần thứ hai: hợp nhất định tuyến}

\subsection{Làm gì}

Thay vì để TII một mình quyết định tác vụ, ta trộn điểm số của TII với điểm số
của đầu RP:
\begin{equation}
\mathrm{fused} = w\, z(\ell^{\mathrm{TII}}) + (1-w)\, z(s^{\mathrm{RP}}),
\qquad w = 0.7 .
\label{eq:route}
\end{equation}

Tác vụ có điểm cao nhất được chọn, rồi \textbf{giao cho DRM/CRM và đầu phân lớp
của HRM-PET xử lý tiếp như bình thường}.

Điểm quan trọng: chúng ta \textbf{không bỏ đi thành phần nào của HRM-PET}. DRM,
CRM, đầu phân lớp chung, cơ chế CRCT, tất cả giữ nguyên. Chỉ bước 1 (TII đơn độc)
được thay bằng bước 1 mới (TII kết hợp đầu RP).

\subsection{Ký hiệu $z$ là gì và tại sao bắt buộc phải có}

$z(\cdot)$ là phép \textbf{chuẩn hóa theo từng ảnh}: trừ đi trung bình rồi chia
cho độ lệch chuẩn, tính riêng trên tập tác vụ đã thấy của chính ảnh đó.

Bỏ bước này đi thì phương pháp hỏng hoàn toàn, và đây là lỗi chúng ta thực sự đã
mắc phải. Logit của TII có biên độ rất rộng, ví dụ trải từ $-20$ đến $+15$, trong
khi điểm số hồi quy ridge của đầu RP có biên độ hẹp, ví dụ từ $-0.3$ đến $+0.8$.
Cộng thẳng hai đại lượng này thì TII át hoàn toàn ở \emph{mọi} giá trị $w$, kể cả
$w = 0.1$.

Khi chưa sửa, độ chính xác định tuyến rơi xuống \textbf{68.24}, thấp hơn cả mức
77.12 của riêng đầu RP. Trộn hai nguồn mà lại tệ hơn từng nguồn riêng lẻ chính là
dấu hiệu của lỗi thang đo. Sau khi chuẩn hóa, cả hai về cùng thang và $w$ mới
thực sự mang nghĩa ``trọng số''.

\subsection{Thêm vào đâu}

\begin{table}[htbp]
\centering
\caption{Vị trí của tầng hợp nhất định tuyến.}
\begin{tabular}{lL}
\toprule
Việc & Vị trí \\
\midrule
Hàm hợp nhất & \texttt{fuse\_routers}, engine dòng 763 \\
Phép chuẩn hóa $z$ & hàm \texttt{standardize} bên trong \texttt{fuse\_routers} \\
Đổi điểm lớp thành điểm tác vụ
  & \texttt{\_task\_scores\_from\_class\_scores}, dòng 719 \\
Tiêm kết quả vào luồng DRM
  & engine dòng 2006, ngay trước \texttt{id\_logits = F.softmax(...)} ở dòng 2015 \\
\bottomrule
\end{tabular}
\end{table}

Vị trí tiêm là chỗ đáng chú ý nhất. Chúng ta chèn đúng vào \textbf{trước} khi
HRM-PET chuẩn hóa điểm định tuyến thành xác suất, nên toàn bộ DRM/CRM phía sau
vẫn chạy bình thường trên đường định tuyến mới. Không dòng nào của HRM-PET phải
sửa.

\noindent Cờ bật:
\texttt{-{}-rp\_route\_fusion\_drm -{}-rp\_route\_fusion\_weight 0.7}

% =====================================================================
\section{Thành phần thứ ba: hợp nhất phân lớp có cổng}

\subsection{Phiên bản đầu tiên và lỗi của nó}

Từ phát hiện 3, ý tưởng hiển nhiên là trộn luôn hai bộ phân lớp với một trọng số
$\beta$ cố định. Kết quả: độ chính xác phân lớp tăng đáng kể trên cả ba bộ dữ
liệu ($+1.40$, $+0.76$, $+1.61$). Nhưng xuất hiện một vấn đề dai dẳng:
\textbf{hàm mất mát trên CIFAR-100 xấu đi $+0.019 \pm 0.003$.}

Con số này nhỏ nhưng phải xử lý, vì độ lệch chuẩn $0.003$ chỉ bằng một phần sáu
giá trị trung bình, nghĩa là nó \textbf{lặp lại y hệt trên mọi seed}. Đó là lỗi
hệ thống, không phải nhiễu ngẫu nhiên. Đối chiếu: chỉ số Backward có độ lệch
chuẩn \emph{lớn hơn} trung bình nhiều lần, và đó mới là nhiễu thật.

\subsection{Truy nguyên}

Trọng số cố định đối xử \textbf{như nhau} với hai loại ảnh có bản chất hoàn toàn
khác: ảnh mà đầu phân lớp chính đã trả lời chắc nịch với xác suất $0.97$, và ảnh
mà nó đang phân vân giữa hai lớp $0.35$ với $0.30$.

Trên CIFAR-100, loại thứ nhất chiếm đa số áp đảo vì mô hình vốn đã đúng gần
$90\,\%$. Dành $30\,\%$ quyền quyết định cho ý kiến thứ hai trên chính những ảnh
đó chỉ \textbf{làm phẳng một đỉnh xác suất vốn đã đúng}. Hàm mất mát phạt ngay
lập tức, trong khi dự đoán cuối cùng chẳng đổi.

Nói ngắn gọn: ta đang trả giá ở $90\,\%$ số ảnh để mua lợi ích ở $10\,\%$ còn lại.

\subsection{Cách sửa: cơ chế cổng theo biên top-2}

Cho $\beta$ thay đổi theo từng ảnh, tỉ lệ nghịch với mức độ chắc chắn của đầu
phân lớp chính. Gọi $p_{(1)}$ và $p_{(2)}$ là hai xác suất cao nhất mà đầu phân
lớp chính đưa ra:
\begin{equation}
\beta_i = \beta \left( 1 - \frac{p_{(1)} - p_{(2)}}{p_{(1)} + p_{(2)}} \right),
\qquad \beta = 0.5 .
\label{eq:gate}
\end{equation}

\begin{luuy}
\textbf{Vì sao chọn biên giữa hai lớp cao nhất, chứ không phải entropy hay xác
suất cao nhất?} Lập luận rất cụ thể: việc trộn chỉ có thể thay đổi dự đoán bằng
cách hất lớp đang đứng đầu xuống. Mà ứng viên khả dĩ nhất để thay thế nó chính là
lớp đứng thứ hai. Vậy đại lượng cần đo là khoảng cách giữa đúng hai lớp đó.
\end{luuy}

\begin{table}[htbp]
\centering
\caption{Cơ chế cổng chạy thực tế.}
\label{tab:gate}
\begin{tabular}{lrrrl}
\toprule
Tình huống & $p_{(1)}$ & $p_{(2)}$ & $\beta_i$ & Hệ quả \\
\midrule
Mô hình rất chắc chắn & 0.90 & 0.02 & $\approx 0.007$ & gần như không đụng tới \\
Mô hình đang phân vân & 0.35 & 0.30 & $\approx 0.28$  & trộn gần hết công suất \\
\bottomrule
\end{tabular}
\end{table}

Đúng như mong muốn: giữ nguyên chỗ mô hình đã đúng, can thiệp mạnh chỗ nó đang
lung lay. Vì cơ chế cổng kéo trọng số hiệu dụng xuống rất nhiều, $\beta$ danh
nghĩa phải nâng từ $0.3$ lên $\mathbf{0.5}$. Chúng ta có thử $0.8$ nhưng lúc đó
trộn quá mạnh và ImageNet-R hỏng cả ba chỉ số Loss, Forgetting, Backward.

Kết quả sau khi thêm cổng: hàm mất mát trên CIFAR-100 đi từ
$\mathbf{+0.019 \pm 0.003}$ về $\mathbf{-0.000 \pm 0.001}$. Suy giảm hệ thống
biến mất hoàn toàn.

\subsection{Bước cuối: đưa về lại thang đo cũ}

Sau khi trộn, kết quả nằm ở thang ``trung bình 0, phương sai 1'', không phải
thang của logit ban đầu. Đưa nó về:
\begin{align}
m_i &= (1 - \beta_i)\, z(\ell_i) + \beta_i\, z(s_i),\\
\hat{\ell}_i &= m_i \cdot \mathrm{std}(\ell_i) + \mathrm{mean}(\ell_i).
\label{eq:affine}
\end{align}

Đây là \textbf{biến đổi affine theo từng ảnh}, tức là đơn điệu, nên \textbf{thứ
tự xếp hạng giữa các lớp không đổi} và mọi chỉ số độ chính xác giữ nguyên đến
từng bit. Nó chỉ đưa thang đo về mức so sánh được với baseline khi tính hàm mất
mát.

\subsection{Thêm vào đâu}

\begin{table}[htbp]
\centering
\caption{Vị trí của tầng hợp nhất phân lớp.}
\begin{tabular}{lL}
\toprule
Việc & Vị trí \\
\midrule
Hàm trộn phân lớp & \texttt{fuse\_class\_scores}, engine dòng 854 \\
Cơ chế cổng & \texttt{\_fusion\_gate}, dòng 824 \\
Chuẩn hóa trên lớp hợp lệ
  & \texttt{\_valid\_moments} dòng 808, \texttt{\_standardize\_valid} dòng 818 \\
Nơi gọi & engine dòng 2178--2179, ngay sau khi đã có \texttt{logits} và trước
          khi tính hàm mất mát \\
\bottomrule
\end{tabular}
\end{table}

\noindent Cờ bật:
\texttt{-{}-rp\_class\_fusion\_weight 0.5 -{}-rp\_class\_fusion\_gate margin}

\begin{luuy}
\textbf{Chi phí bằng không.} Điểm số của đầu RP đã được tính ở tầng định tuyến
rồi; tầng này chỉ dùng lại. Toàn bộ phần tăng thêm của hệ thống là đúng
\emph{một} lượt chạy mạng cho đầu RP.
\end{luuy}

\begin{figure}[htbp]
\centering
\resizebox{\textwidth}{!}{%
\begin{tikzpicture}[
  font=\footnotesize, >={Stealth[length=2mm]},
  box/.style={draw, minimum height=8mm, minimum width=1.6cm, inner xsep=4pt, align=center},
  hl/.style={box, very thick, fill=black!8}
]
\node[box] (in) at (0.0, 0.0)    {Ảnh vào};
\node[box] (tii) at (2.7, 0.85)  {TII};
\node[box] (rp) at (2.7, -0.85)  {Đầu RP};
\node[hl]  (f1) at (5.9, 0.0)    {Tầng 1\\Hợp nhất định tuyến};
\node[box] (lo) at (8.9, 0.0)    {Chọn LoRA};
\node[box] (dr) at (11.3, 0.0)   {DRM / CRM};
\node[box] (cl) at (13.9, 0.0)   {Đầu phân lớp};
\node[hl]  (f2) at (13.9, -2.1)  {Tầng 2\\Hợp nhất phân lớp};
\node[box] (out) at (13.9, -3.7) {Dự đoán};

\draw[->] (in) -- (tii);
\draw[->] (in) -- (rp);
\draw[->] (tii) -- (f1);
\draw[->] (rp) -- (f1);
\draw[->] (f1) -- (lo);
\draw[->] (lo) -- (dr);
\draw[->] (dr) -- (cl);
\draw[->] (cl) -- (f2);
\draw[->] (f2) -- (out);
\draw[->, dashed] (rp.south) |- node[pos=0.75, above, font=\scriptsize]
      {điểm số theo lớp của đầu RP} (f2.west);
\end{tikzpicture}}
\caption{Quy trình sau khi bổ sung. Hai khối viền đậm là thành phần được thêm.
Đầu RP được tính \emph{một lần duy nhất} và cung cấp tín hiệu cho cả hai tầng:
điểm số theo tác vụ cho tầng thứ nhất, điểm số theo lớp cho tầng thứ hai (đường
nét đứt). So sánh với Hình \ref{fig:goc} để thấy chính xác phần được thêm.}
\label{fig:moi}
\end{figure}

% =====================================================================
\section{Bốn cạm bẫy đã gặp}

Ghi lại để người sau không mất công lặp lại.

\subsection{Lệch thang đo khi trộn}

Dấu hiệu nhận biết: \textbf{kết quả trộn tệ hơn từng nguồn riêng lẻ}. Nếu gặp lại
triệu chứng này ở bất kỳ phép trộn nào, hãy nghi ngờ thang đo trước tiên.

Cách kiểm chứng đã dùng: đặt $w = 1.0$ phải tái lập baseline \textbf{chính xác
đến từng chữ số}. Chúng ta xác nhận được $77.7914$, $74.0191$, $1.2305$ trên ba
bộ dữ liệu, không sai khác ở chữ số thập phân thứ tư.

\subsection{Lỗi chuẩn hóa ảnh đầu vào -- lỗi thật của repo}

Hàm \texttt{build\_transform} trong \texttt{datasets.py} dòng 347 gọi
\texttt{ToTensor()} mà \textbf{không} gọi \texttt{Normalize()}. Ảnh do đó nằm
trong khoảng $[0,1]$, trong khi checkpoint ViT-B/16 của ImageNet-21K mong đợi
khoảng $[-1,1]$.

HRM-PET miễn nhiễm với lỗi này vì nó \emph{huấn luyện} trên chính những đặc trưng
bị lệch đó. Nhưng một đầu phân lớp không huấn luyện như đầu RP thì không miễn
nhiễm. Sửa lại đúng khoảng giá trị đem lại \textbf{$+2.7$ điểm} trên đặc trưng
đóng băng của CUB-200.

Cảnh báo kèm theo: \texttt{build\_cifar\_transform} (dòng 397) \emph{có} gọi
\texttt{Normalize} với thống kê CIFAR. Chồng thêm phép chuẩn hóa lên đó làm CIFAR
sập xuống 50.42. Vì vậy có hai chế độ riêng, \texttt{half} và
\texttt{cifar\_half}.

\subsection{Lệnh chạy nền tự tìm thấy chính nó}

Vòng lặp chờ GPU ban đầu viết là \texttt{pgrep -f "\#\#\# SURVEY"} đặt ngay trong
chuỗi \texttt{bash -c}. Nhưng \texttt{pgrep -f} so khớp với \textbf{toàn bộ dòng
lệnh} của mọi tiến trình, và dòng lệnh của chính tiến trình đó chứa nguyên văn
chuỗi đang tìm. Nó tự tìm thấy mình và chờ vĩnh viễn.

Cách sửa: đưa mẫu tìm kiếm vào một file script riêng để nó không còn nằm trên
dòng lệnh của bên gọi. Xem \texttt{training\_scripts/wait\_for\_gpu.sh}.

\subsection{\texttt{grep} tự đọc và ghi cùng một file}

Có lần lệnh viết dạng \texttt{nohup bash -c '\ldots; grep \ldots{} \textasciitilde/cost.log'
> \textasciitilde/cost.log}. Lệnh \texttt{grep} vừa đọc vừa ghi cùng một file, sinh
vòng lặp vô hạn và \textbf{lấp đầy ổ đĩa đến $99\,\%$}. Quy tắc rút ra: không bao
giờ grep chính file mà lệnh đang ghi vào.

% =====================================================================
\section{Kết quả và cách đọc chúng}

\subsection{Sáu chỉ số nghĩa là gì}

\begin{table}[htbp]
\centering
\caption{Ý nghĩa các chỉ số.}
\begin{tabular}{llc}
\toprule
Chỉ số & Nghĩa & Chiều tốt \\
\midrule
Acc@task   & tỉ lệ ảnh được gán đúng \emph{tác vụ} & cao \\
Acc@1      & tỉ lệ ảnh được gán đúng \emph{lớp} -- quan trọng nhất & cao \\
Acc@5      & tỉ lệ ảnh có lớp đúng nằm trong 5 lớp dẫn đầu & cao \\
Loss       & hàm mất mát, đo cả độ chắc chắn chứ không chỉ đúng/sai & thấp \\
Forgetting & mức tụt so với đỉnh của chính tác vụ đó & thấp \\
Backward   & chênh lệch giữa độ chính xác cuối và ngay sau khi học & cao \\
\bottomrule
\end{tabular}
\end{table}

Hai chỉ số cuối cần đọc cẩn thận: chúng so một tác vụ \textbf{với chính nó ở thời
điểm trước}, chứ không so giữa các phương pháp. Đặc điểm này sẽ quay lại ở Mục
\ref{sec:ramp}.

\subsection{Bảng kết quả, bốn seed}

\begin{table}[htbp]
\centering
\caption{Kết quả trên ba bộ dữ liệu, trung bình $\pm$ độ lệch chuẩn trên bốn seed
42, 43, 44, 45.}
\label{tab:main}
\begin{tabular}{@{}p{0.4cm}lRRR@{}}
\toprule
\multicolumn{2}{@{}l}{Bộ dữ liệu và chỉ số} & HRM-PET & Đề xuất & Thay đổi \\
\midrule
\multicolumn{5}{l}{\emph{Split-ImageNet-R}}\\
& Acc@task   & $77.58 \pm 0.33$  & $79.85 \pm 0.18$  & $+2.264 \pm 0.158$ \\
& Acc@1      & $73.94 \pm 0.48$  & $75.18 \pm 0.24$  & $+1.236 \pm 0.416$ \\
& Acc@5      & $86.71 \pm 0.21$  & $88.02 \pm 0.17$  & $+1.312 \pm 0.148$ \\
& Loss       & $1.23 \pm 0.01$   & $1.19 \pm 0.01$   & $-0.040 \pm 0.002$ \\
& Forgetting & $3.36 \pm 0.56$   & $3.30 \pm 0.35$   & $-0.056 \pm 0.334$ \\
& Backward   & $-3.16 \pm 0.61$  & $-3.20 \pm 0.42$  & $-0.047 \pm 0.282$ \\
\addlinespace
\multicolumn{5}{l}{\emph{Split-CIFAR100}}\\
& Acc@task   & $89.77 \pm 0.11$  & $90.87 \pm 0.09$  & $+1.103 \pm 0.066$ \\
& Acc@1      & $89.78 \pm 0.03$  & $90.47 \pm 0.05$  & $+0.690 \pm 0.063$ \\
& Acc@5      & $98.65 \pm 0.03$  & $98.86 \pm 0.07$  & $+0.217 \pm 0.057$ \\
& Loss       & $0.39 \pm 0.00$   & $0.39 \pm 0.00$   & $-0.000 \pm 0.001$ \\
& Forgetting & $3.89 \pm 0.08$   & $3.74 \pm 0.12$   & $-0.156 \pm 0.098$ \\
& Backward   & $-3.87 \pm 0.07$  & $-3.73 \pm 0.11$  & $+0.142 \pm 0.097$ \\
\addlinespace
\multicolumn{5}{l}{\emph{Split-CUB200}}\\
& Acc@task   & $93.12 \pm 0.14$  & $94.34 \pm 0.10$  & $+1.226 \pm 0.149$ \\
& Acc@1      & $86.38 \pm 0.25$  & $87.88 \pm 0.27$  & $+1.495 \pm 0.353$ \\
& Acc@5      & $97.04 \pm 0.10$  & $97.44 \pm 0.13$  & $+0.400 \pm 0.103$ \\
& Loss       & $0.57 \pm 0.00$   & $0.53 \pm 0.01$   & $-0.036 \pm 0.004$ \\
& Forgetting & $2.54 \pm 0.35$   & $2.18 \pm 0.12$   & $-0.359 \pm 0.325$ \\
& Backward   & $-2.42 \pm 0.39$  & $-2.14 \pm 0.13$  & $+0.279 \pm 0.391$ \\
\bottomrule
\end{tabular}
\end{table}

\subsection{Cách đọc cột ``Thay đổi''}

Cột này \textbf{không phải} hiệu của hai cột bên trái. Nó được tính theo cặp: với
từng seed, lấy kết quả đề xuất trừ kết quả baseline, rồi mới lấy trung bình của
bốn hiệu đó.

Khác biệt nằm ở độ lệch chuẩn. Có seed dễ có seed khó; nếu lấy hiệu của hai trung
bình, độ biến thiên do độ khó của seed sẽ trộn lẫn vào. Tính theo cặp thì độ khó
của seed triệt tiêu, và độ lệch chuẩn còn lại \textbf{chỉ phản ánh độ ổn định của
chính mức cải thiện}. Quy tắc đọc:

\begin{itemize}
\item \textbf{Độ lệch chuẩn nhỏ hơn trung bình nhiều lần} $\rightarrow$ cải thiện
      thật, lặp lại được. Ví dụ Acc@task trên CIFAR-100: $+1.103 \pm 0.066$, tỉ
      lệ 17 lần.
\item \textbf{Độ lệch chuẩn lớn hơn trung bình} $\rightarrow$ không kết luận được
      gì. Ví dụ Backward trên ImageNet-R: $-0.047 \pm 0.282$, độ lệch gấp sáu lần
      trị tuyệt đối của trung bình. Đây là ô duy nhất không dương trong toàn bộ
      18 phép so sánh, \textbf{và nó không phải suy giảm, nó là nhiễu}.
\end{itemize}

\subsection{Chi phí}

\begin{table}[htbp]
\centering
\caption{Số lượt gọi LoRA trung bình cho mỗi ảnh.}
\begin{tabular}{lrrr}
\toprule
Bộ dữ liệu & HRM-PET & Đề xuất & Tăng thêm \\
\midrule
ImageNet-R & 1.83 & 2.85 & $+1.02$ \\
CIFAR-100  & 1.40 & 2.41 & $+1.01$ \\
CUB-200    & 1.60 & 2.60 & $+1.00$ \\
\bottomrule
\end{tabular}
\end{table}

Mức tăng đúng bằng \emph{một} lượt, chính là lượt chạy cho đầu RP. Tầng hợp nhất
phân lớp không tốn thêm gì. Bộ nhớ tăng thêm gồm một ma trận Gram, một ma trận
tương quan lớp và một seed.

% =====================================================================
\section{Cách chạy lại toàn bộ}

\subsection{Cấu hình đã chốt}

\begin{table}[htbp]
\centering
\caption{Cấu hình cuối cùng.}
\begin{tabular}{lll}
\toprule
Biến & Giá trị & Ý nghĩa \\
\midrule
\texttt{RP\_SOURCE}   & \texttt{lora}   & đặc trưng lấy từ LoRA, không phải mạng đóng băng \\
\texttt{RP\_INORM}    & \texttt{none}   & không đổi khoảng giá trị đầu vào \\
\texttt{RP\_FUSE}     & \texttt{1}      & bật hợp nhất định tuyến \\
\texttt{RP\_FUSE\_DRM}& \texttt{1}      & tiêm kết quả vào luồng DRM/CRM \\
\texttt{RP\_FUSE\_W}  & \texttt{0.7}    & trọng số trên TII \\
\texttt{RP\_CLS\_W}   & \texttt{0.5}    & trọng số danh nghĩa của hợp nhất phân lớp \\
\texttt{RP\_CLS\_GATE}& \texttt{margin} & cơ chế cổng theo biên top-2 \\
\texttt{RP\_RAMP}     & \texttt{0.0}    & tắt cơ chế tăng dần, xem Mục \ref{sec:ramp} \\
\texttt{CALIBRATE}    & \texttt{0}      & không hiệu chỉnh nhiệt độ \\
\texttt{RP\_DIM}      & \texttt{10000}  & số chiều phép chiếu ngẫu nhiên \\
\texttt{RP\_LAMBDA}   & \texttt{10000}  & hệ số điều chuẩn của hồi quy ridge \\
\bottomrule
\end{tabular}
\end{table}

\subsection{Chạy đánh giá cho một seed}

\begin{footnotesize}
\begin{verbatim}
cd ~/Documents/truongnguyen/Hybrid_ReMatching && nohup bash -c \
  'bash training_scripts/wait_for_gpu.sh 12000 120; \
   RP_SOURCE=lora RP_INORM=none RP_FUSE=1 RP_FUSE_DRM=1 RP_FUSE_W=0.7 \
   RP_CLS_W=0.5 RP_CLS_GATE=margin CALIBRATE=0 RP_DIM=10000 RP_LAMBDA=10000 \
   bash training_scripts/rollout_hybrid_4090.sh 42' \
  > ~/run.log 2>&1 & tail -F --pid=$! ~/run.log
\end{verbatim}
\end{footnotesize}

\subsection{Lấy bảng kết quả}

\begin{footnotesize}
\begin{verbatim}
cd ~/Documents/truongnguyen/Hybrid_ReMatching && \
  ./.venv/bin/python training_scripts/summarize_hybrid_multiseed.py \
  ~/Documents/truongnguyen/hrm-pet-output \
  "_eval_rp_lora_*f1d1w0p7*cw0p5sh1p0m1gmargin.log"
\end{verbatim}
\end{footnotesize}

\subsection{File nào làm gì}

\begin{table}[htbp]
\centering
\caption{Bản đồ mã nguồn.}
\begin{tabular}{lL}
\toprule
File & Vai trò \\
\midrule
\texttt{engines/random\_projection\_head.py}
  & toàn bộ đầu RP: chiếu, tích lũy, giải ridge \\
\texttt{engines/hrm\_lora\_wtp\_and\_tap\_engine.py}
  & hai tầng hợp nhất và điểm tiêm vào HRM-PET \\
\texttt{trainers/lora\_trainer.py}
  & gọi tích lũy và giải đầu RP sau mỗi tác vụ \\
\texttt{configs/imr\_lora.py}
  & khai báo cờ cho ImageNet-R và CUB-200 \\
\texttt{configs/cifar100\_lora.py}
  & khai báo cờ cho CIFAR-100 \\
\texttt{training\_scripts/eval\_rp\_head\_any\_4090.sh}
  & chạy đánh giá cho một bộ dữ liệu \\
\texttt{training\_scripts/rollout\_hybrid\_4090.sh}
  & chạy cả ba bộ dữ liệu cho một seed \\
\texttt{training\_scripts/summarize\_hybrid\_multiseed.py}
  & bảng trung bình $\pm$ độ lệch chuẩn \\
\texttt{training\_scripts/incremental\_accuracy.py}
  & độ chính xác theo từng giai đoạn \\
\texttt{training\_scripts/wait\_for\_gpu.sh}
  & chờ GPU trống rồi mới chạy \\
\bottomrule
\end{tabular}
\end{table}

% =====================================================================
\section{Những hướng đã thử và đã đóng}

Phần này quan trọng không kém phần thành công. Ghi lại để không ai mất công làm
lại.

\subsection{Re-matching vét cạn}

Thử mọi LoRA cho mọi ảnh rồi chọn kết quả tốt nhất. Kết quả \textbf{không vượt
được baseline}, dù tốn gấp 10 lần. Nguyên nhân: đầu phân lớp dùng chung khiến
LoRA của tác vụ khác vẫn có thể tạo ra điểm số cao cho lớp không thuộc nó, hiện
tượng ``chiếm quyền'' giữa các tác vụ.

\subsection{Tái chấm điểm bằng Gaussian hoặc Mahalanobis}

Mô hình hóa đặc trưng 768 chiều của mỗi tác vụ bằng một Gaussian rồi định tuyến
theo khoảng cách Mahalanobis. Định tuyến rơi xuống \textbf{84.3 so với 93.21}.
Thử cả hiệp phương sai riêng lẫn hiệp phương sai chung, đều hỏng.

Nguyên nhân cấu trúc: LoRA gắn trên mạng đóng băng làm đặc trưng thay đổi
\emph{quá ít} để tạo được tín hiệu phân biệt trong ngoài phân bố.

\subsection{Tăng hạng LoRA từ 8 lên 16}

Độ chính xác gần như không đổi (86.46 so với 86.53). Kết luận: \textbf{đây không
phải vấn đề dung lượng mô hình}, nên tăng tham số không giúp gì.

\subsection{Đầu RP dùng một mình, bỏ hẳn định tuyến}

Thắng trên CUB-200 (87.96 so với 86.53) nhưng thua nặng trên ImageNet-R (70.08 so
với 74.02). Lý do rất rõ: CUB-200 là ảnh chim, mà ImageNet-21K vốn giàu lớp chim,
nên đặc trưng đóng băng đã đủ tốt. ImageNet-R là ảnh phong cách hóa, xa phân bố
tiền huấn luyện, nên \textbf{bắt buộc phải có thích nghi theo tác vụ}, và đó
chính là thứ các LoRA của HRM-PET cung cấp.

Bài học: hai bên bổ trợ nhau, không thay thế nhau.

\subsection{Thống kê theo tầng, lấy từ bài báo PMI-CFS}

Đọc kỹ cả phụ lục, chỉ có \emph{một} thành phần khả dụng: thống kê kích hoạt theo
tầng làm bộ định tuyến thứ ba. Đã cài ở
\texttt{engines/layer\_stat\_router.py} nhưng chưa chạy hoàn chỉnh; hướng này đã
dừng.

Ba thành phần còn lại không áp dụng được. PMI cần sinh ảnh giả, mà HRM-PET phát
lại ở không gian đặc trưng nên không cần. CFS đã có sẵn trong repo và cho kết quả
\emph{thấp hơn} baseline. Phép chiếu ngữ nghĩa cần bộ mã hóa văn bản CLIP mà kiến
trúc của ta không có.

\subsection{Trọng số tăng dần theo giai đoạn -- đã cài, đã đo, đã loại}
\label{sec:ramp}

Đây là ví dụ đáng học nhất về việc \textbf{một cải tiến có vẻ hoàn hảo lại phải
bị loại}.

Xuất phát từ một quan sát đúng: Forgetting và Backward gần như không nhúc nhích.
Lý do là cả hai đều so một tác vụ với đỉnh của chính nó, nên một hiệu chỉnh áp
dụng đồng đều ở mọi giai đoạn sẽ \textbf{triệt tiêu khỏi cả hai}.

Giải pháp thử nghiệm: cho trọng số trộn tăng dần theo số tác vụ đã thấy, dạng
$(\text{số tác vụ đã thấy}/T)^{\gamma}$. Vì hệ số này đạt $1.0$ ở giai đoạn cuối
nên hàng kết quả cuối cùng \textbf{không đổi một chữ số nào}, chỉ hai chỉ số ma
trận mới thay đổi.

Và nó chạy đúng như thiết kế. Với $\gamma = 2$, Forgetting cải thiện lên $-0.415$,
$-0.433$, $-0.462$ và Backward lên $+0.357$, $+0.411$, $+0.462$ trên ba bộ dữ
liệu.

\begin{luuy}
\textbf{Nhưng chính điều đó tố cáo nó.} Nếu trạng thái cuối không đổi mà hai chỉ
số so-với-đỉnh lại cải thiện, thì mức cải thiện đó chỉ có thể đến từ việc mô hình
\emph{kém đi ở các giai đoạn giữa}. Đo độ chính xác trung bình qua cả mười giai
đoạn xác nhận đúng như vậy: trên CUB-200 nó tụt từ $88.35$ xuống $88.16$ rồi
$87.91$. Cải thiện chỉ số ghi nhớ bằng cách hạ thấp đường đi thì không phải cải
tiến, và phản biện sẽ chỉ ra ngay.
\end{luuy}

Cơ chế vẫn còn trong mã (\texttt{-{}-rp\_fusion\_ramp}) làm dòng ablation, nhưng
\textbf{mặc định tắt}. Hai điều có giá trị vẫn thu được:

\begin{enumerate}
\item Hợp nhất phân lớp \emph{thực sự} gây hại ở giai đoạn đầu: ImageNet-R chỉ
      đạt $86.0$ so với $87.5$ của baseline ở tác vụ đầu tiên, đúng như dự đoán
      rằng ước lượng ma trận Gram cần đủ số lớp mới đáng tin.
\item Nhưng cùng đầu RP đó lại là \emph{bộ định tuyến mạnh ngay từ đầu}. Nên giảm
      trọng số định tuyến ở giai đoạn sớm chỉ làm mất độ chính xác định tuyến mà
      không được gì bù lại.
\end{enumerate}

\subsection{Bài học về phương pháp}

Trước khi chạy thí nghiệm ramp, chúng ta \textbf{ghi trước ba điều kiện nghiệm
thu}. Kết quả: một điều kiện đạt, một điều kiện trượt, và một điều kiện hóa ra
không phải phép thử (chỉ số Backward trên seed 42 vốn đã dương sẵn nên không kiểm
chứng được gì).

Nếu không ghi trước, thí nghiệm này sẽ trông như một thắng lợi sạch trên mọi chỉ
số của bảng tổng kết. Việc ghi trước là thứ đã bắt được nó.

% =====================================================================
\section*{Trạng thái hiện tại}
\addcontentsline{toc}{section}{Trạng thái hiện tại}

\begin{itemize}
\item Cấu hình đã chốt, đủ bốn seed trên cả ba bộ dữ liệu, \textbf{17 trong 18
      phép so sánh cải thiện}.
\item Báo cáo kỹ thuật ở \texttt{reports/bao\_cao\_fusion.tex}.
\item Việc còn lại đáng làm nhất: bổ sung \textbf{ImageNet-A} làm bộ dữ liệu thứ
      tư. Phần mã đã sẵn sàng -- lớp nạp dữ liệu ở
      \texttt{continual\_datasets/continual\_datasets.py} dòng 712 và cấu hình ở
      \texttt{configs/ima\_lora.py}. Thứ duy nhất còn thiếu là bản thân dữ liệu,
      phải tải thủ công về \texttt{datasets/imagenet-a/}.
\end{itemize}

\end{document}
